% !TeX root = ../../Book2.tex
% 纲要标题：### 9.1 交替多线性映射
% 纲要位置：计划纲要.md，第 1068 行。

\section{交替多线性映射}
\label{b2:sec:0901}

面积在交换两条边时改变符号，在两条边重合时变为零。这两种性质在熟悉的实数情形中相互联系，但在一般系数环上并不等价。外幂所表达的是“有两个输入相同时为零”的交替性。本章始终设 \(R\) 为含单位元的交换环，所有张量积和多线性映射均在 \(R\) 上；需要域或额外特征条件时另行说明。

% 纲要要点（供目录核对）：
%
% * 交替性与反对称性的区别；双线性矩阵的对角线条件
% * 特征二与目标模中的二挠现象
% * 外幂的泛性质；先用张量积表达多线性，再按重复输入取商
% * 外幂不保持一般单射的理想反例，完整证明放在正文
%

\subsection{交替性与反对称性}
\label{b2:subsec:0901:01}

\begin{definition}\label{b2:def:alternating-multilinear}
设 \(R\) 为交换环，\(M,L\) 为 \(R\) 模，\(p\ge1\)。对逐变量 \(R\) 线性的映射 \(b:M^p\to L\)，若任意两个输入相同便有
\[
b(m_1,\ldots,m_p)=0,
\]
就称 \(b\) 为\term[jiaoti-duoxianxing-yingshe]{交替多线性映射}[alternating multilinear map]。若对任意两个位置，交换输入后 \(b\) 的值都变为原值的负元，就称 \(b\) 为\term[fanduichen-duoxianxing-yingshe]{反对称多线性映射}[skew-symmetric multilinear map]。\(p=1\) 时，两种条件均为空条件，此时这两类映射就是线性映射。
\end{definition}

\begin{proposition}\label{b2:prop:alternating-implies-skew}
设 \(R\) 为交换环，\(M,L\) 为 \(R\) 模，\(p\ge1\)。从 \(M^p\) 到 \(L\) 的交替多线性映射必反对称。反过来，若目标模 \(L\) 中 \(2\ell=0\) 蕴含 \(\ell=0\)，则从 \(M^p\) 到 \(L\) 的反对称多线性映射也交替。特别地，当 \(2\) 在 \(R\) 中可逆时，两种条件等价。
\end{proposition}
\begin{proof}
\(p=1\) 时两个条件均为空条件。以下设 \(p\ge2\)。
固定除第 \(i,j\) 个位置以外的所有输入，把这两个位置都取为 \(x+y\)。交替性使函数值为零；按这两个变量展开，两个含 \((x,x)\)、\((y,y)\) 的项分别为零，故剩下的两项满足
\[
b(\ldots,x,\ldots,y,\ldots)
+b(\ldots,y,\ldots,x,\ldots)=0.
\]
这就是反对称性。

若反对称映射在两个位置取相同的 \(x\)，交换它们不改变输入，却使值变号，所以这个值 \(\ell\) 满足 \(2\ell=0\)。所给假设使 \(\ell=0\)，即交替。当 \(2\) 可逆时，乘以 \(2^{-1}\) 就能从 \(2\ell=0\) 推出 \(\ell=0\)。
\end{proof}

将置换写成对换的复合，交替性还给出
\begin{equation}\label{b2:eq:alternating-permutation}
b(m_{\sigma(1)},\ldots,m_{\sigma(p)})
=\operatorname{sgn}(\sigma)b(m_1,\ldots,m_p)
\qquad(\sigma\in S_p).
\end{equation}
所用的符号同态及其在换位上取值为 \(-1\) 的性质，已由第一卷定理3.2.3证明。符号 \(\operatorname{sgn}(\sigma)\) 只由置换决定，因此表达不依赖所选的对换分解；公式在 \(p!\) 不可逆的环上也成立。

\ProseParagraphBreak
例如，任意交换环上的
\[
b\bigl((x_1,x_2),(y_1,y_2)\bigr)=x_1y_2-x_2y_1
\]
是交替双线性映射。其交替性应直接由 \(x_1x_2-x_2x_1=0\) 核对；不能先用“交换变号”，再在可能存在二挠的环上消去一个 \(2\)。

\subsection{特征二与二挠现象}
\label{b2:subsec:0901:02}

在域 \(k=\mathbb F_2\) 上，双线性映射 \(b:k\times k\to k\)、\(b(x,y)=xy\) 满足
\[
b(y,x)=b(x,y)=-b(x,y),
\]
所以反对称，却有 \(b(1,1)=1\)，因而不交替。在特征二中，“交换变号”变成“交换不变”，完全没有强迫对角线取值为零。

\ProseParagraphBreak
这一区别不只由环的特征决定。取 \(R=\mathbb Z\)、\(M=\mathbb Z\)、\(L=\mathbb Z/2\mathbb Z\)，公式 \(b(a,c)=\overline{ac}\) 同样反对称而不交替。环本身特征零，目标模却有二挠。因此，使用\cref{b2:prop:alternating-implies-skew}从反对称性推出交替性时，要检查的是实际取值所在的目标模能否消去 \(2\)，而非仅仅检查底环的特征。

\ProseParagraphBreak
对于交换环 \(R\) 上基为 \(e_1,\ldots,e_n\) 的有限自由模 \(M\)，一个双线性形式 \(b:M\times M\rightarrow R\) 由矩阵 \(A=(a_{ij})\)、\(a_{ij}=b(e_i,e_j)\) 确定。反对称性等价于 \(A^{\mathsf t}=-A\)，而交替性还要求对角线为零，即
\[
a_{ii}=0,\qquad a_{ij}=-a_{ji}.
\]
必要性由交替性及\cref{b2:prop:alternating-implies-skew}得到。反过来，对 \(x=\sum_i x_ie_i\)，展开
\[
b(x,x)=\sum_i a_{ii}x_i^2+
\sum_{i<j}(a_{ij}+a_{ji})x_ix_j=0.
\]
所以条件充分。特征二中，这种矩阵是对称且对角线为零的矩阵；一般对称矩阵并不足够。后续讨论交替双线性型时，仍以对角线取值为零作为交替性的定义。

\subsection{外幂的泛性质}
\label{b2:subsec:0901:03}

前面的区别决定了外幂应当施加哪些关系。如果在二重张量积中只把 \(x\otimes y+y\otimes x\) 置为零，特征二时这只强制交换不变，并不强制 \(x\otimes x=0\)。例如 \(M=\mathbb F_2 e\) 时，所有这样的反对称关系本来就是零，非零的 \(e\otimes e\) 会被保留下来。因此，要表达交替性，必须直接把重复输入所对应的张量置为零。

\begin{definition}\label{b2:def:exterior-power}
设 \(R\) 为交换环，\(M\) 为 \(R\) 模。对 \(p\ge2\)，在 \(p\) 重张量积 \(M^{\otimes p}\) 中，令 \(J_p\) 为所有具有两个相同因子的纯张量生成的子模。将商模
\[
\bigwedge^p M\coloneqq M^{\otimes p}/J_p,
\]
称为 \(M\) 的 \(p\) 次\term[waimi]{外幂}[exterior power]，并把 \(m_1\otimes\cdots\otimes m_p\) 的像记为 \(m_1\wedge\cdots\wedge m_p\)。约定 \(\bigwedge^0M=R\)、\(\bigwedge^1M=M\)，空楔积为 \(1\in R\)。
\end{definition}
\symidx{\(\bigwedge^p M\)：模的第 \(p\) 次外幂}
\symidx{\(m_1\wedge\cdots\wedge m_p\)：外幂中的纯楔积}

多重张量积已经把 \(p\) 个输入的多线性关系转化为模中的线性关系；再模掉 \(J_p\)，就是在这个模中加入“重复输入为零”的关系。这与第三章按关系取模的商、以及第八章通过张量商施加代数关系的构造相同。所得映射 \((m_1,\ldots,m_p)\mapsto m_1\wedge\cdots\wedge m_p\) 多线性且交替，所有这些纯楔积生成外幂。

\begin{proposition}[外幂的泛性质]\label{b2:prop:exterior-alternating-universal}
设 \(R\) 为交换环，\(M,L\) 为 \(R\) 模，\(p\ge1\)。对任意交替 \(p\) 线性映射 \(b:M^p\to L\)，存在唯一 \(R\) 线性映射
\[
\widetilde b:\bigwedge^pM\longrightarrow L,
\qquad \widetilde b(m_1\wedge\cdots\wedge m_p)=b(m_1,\ldots,m_p).
\]
反过来，每个这样的线性映射都给出交替 \(p\) 线性映射。具有这一性质的对象连同其指定映射，在唯一的相容同构意义下确定。
\end{proposition}
\begin{proof}
\(p=1\) 时，\(\bigwedge^1M=M\)，结论就是线性映射本身。以下设 \(p\ge2\)。由\cref{b2:thm:multiple-tensor-universal}，多线性映射 \(b\) 唯一延伸为 \(M^{\otimes p}\to L\)。交替性恰好说明 \(J_p\) 的每个生成元都被送到零，所以它唯一地通过商模分解。反过来，商映射本身交替，与线性映射复合仍然交替。唯一性来自纯楔积的生成性。

若另一个对象也具有同样的性质，两次应用泛性质得到双向线性映射；它们的复合在每个指定纯楔积上为恒等，故处处为恒等，证明唯一同构。\(p=0\) 时把零变量的值看成指定的 \(\ell\in L\)，它唯一对应 \(R\to L\)、\(r\mapsto r\ell\)。
\end{proof}

因此，要定义一个从 \(\bigwedge^pM\) 到 \(L\) 的线性映射，可以先给出 \(M^p\) 上的公式，检查它对各变量线性且在重复输入时为零，再由泛性质取得唯一的线性映射。后面读取外幂坐标、定义楔积及收缩算子时，使用的都是这一方法。

\ProseParagraphBreak
设 \(f:M\to N\) 线性。交替公式
\[
(m_1,\ldots,m_p)\longmapsto
f(m_1)\wedge\cdots\wedge f(m_p)
\]
诱导 \(\bigwedge^p f:\bigwedge^pM\to\bigwedge^pN\)。在纯楔积上计算得到
\begin{equation}\label{b2:eq:exterior-functoriality}
\bigwedge^p(gf)=\left(\bigwedge^p g\right)\left(\bigwedge^p f\right),
\qquad \bigwedge^p\operatorname{id}_M=\operatorname{id}_{\bigwedge^pM}.
\end{equation}
于是每个外幂都是一个函子；\(p=0\) 时规定所有 \(\bigwedge^0 f\) 都为 \(\operatorname{id}_R\)。与张量积相似，纯楔积的表达也不唯一，外幂的一般元素是有限和，未必能写成一个纯楔积。

\ProseParagraphBreak
函子化保证每个线性映射都诱导外幂上的映射，却不保证它保留单射。由张量积取商得到外幂，并没有消除第七章讨论过的这种障碍。

\begin{proposition}\label{b2:prop:exterior-injection-failure}
设 \(k\) 为域，\(R=k[x,y]\)，\(I=(x,y)\) 为 \(R\) 的理想，并把 \(I\) 视为 \(R\) 模。包含映射 \(i:I\hookrightarrow R\) 是单射，但诱导映射
\[
\bigwedge^2i:\bigwedge^2_R I\longrightarrow\bigwedge^2_R R
\]
不是单射：\(x\wedge y\ne0\) 于 \(\bigwedge^2_R I\)，而其像为零。
\end{proposition}
\begin{proof}
为了检测 \(x\wedge y\)，只需构造一个在 \((x,y)\) 处不为零的交替双线性映射。把 \(k\) 视为经 \(R\rightarrow R/I\cong k\)、\(r\mapsto r(0,0)\) 取得标量作用的 \(R\) 模。任意 \(f,g\in I\) 唯一地写成
\[
f=ax+by+f_{\ge2},\qquad g=cx+dy+g_{\ge2},
\qquad a,b,c,d\in k,
\]
其中余下的项全次数至少为二。定义 \(\beta(f,g)=ad-bc\)。读取一次项系数是可加的，所以 \(\beta\) 对两个变量均可加。对任意 \(r\in R\)，\(rf\) 的一次项恰为 \(r(0,0)(ax+by)\)，因此
\[
\beta(rf,g)=r(0,0)\beta(f,g),
\qquad \beta(f,rg)=r(0,0)\beta(f,g).
\]
这就是以 \(k\cong R/I\) 为目标的 \(R\) 双线性。又 \(\beta(f,f)=ab-ba=0\)，所以 \(\beta\) 交替。由\cref{b2:prop:exterior-alternating-universal}，它诱导线性映射 \(\widetilde\beta:\bigwedge^2_RI\rightarrow k\)，满足 \(\widetilde\beta(x\wedge y)=\beta(x,y)=1\)。于是 \(x\wedge y\ne0\)。

另一方面，\(R\) 作为自身的模由 \(1\) 生成；任意纯楔积 \(r\wedge s=rs(1\wedge1)=0\)。纯楔积生成二次外幂，所以 \(\bigwedge^2_R R=0\)。故 \(\bigwedge^2i\) 将上述非零元素送到零，不是单射。
\end{proof}
